\documentclass[12pt]{article}
\usepackage[utf8]{inputenc}
\usepackage[T1]{fontenc}
\usepackage[margin=1in]{geometry}
\usepackage{setspace}
\usepackage{xcolor}
\usepackage[colorlinks=true, urlcolor=blue!60!black, linkcolor=blue!60!black]{hyperref}
\setlength{\parindent}{0em}
\setlength{\parskip}{0.6em}
\renewcommand{\baselinestretch}{1.0}
\begin{document}
\begin{center}
{\LARGE \textit{Garcia v. San Antonio Metropolitan Transit Authority}, 469 U.S. 528 (1985)\par}
\vspace{0.5em}
{\large Case Brief --- Two-Page Edition\par}
\end{center}

\section{Memory Jogger}
Overruling \textit{National League of Cities}, the Court held 5--4 that a city transit authority must obey the FLSA; the States' protection lies in the political process (\href{https://supreme.justia.com/cases/federal/us/469/528/}{Justia}).

\section{Detailed Case Facts}
The FLSA's wage and overtime rules reached nearly all state and local employees by 1974. After \textit{National League of Cities} (1976) immunized ``traditional governmental functions,'' San Antonio's public transit system stopped following the overtime rules, though it relied heavily on federal transit grants.

\section{Procedural History}
The Labor Department found SAMTA covered; SAMTA sued for a declaratory judgment, and the district court twice held transit a traditional function. The Supreme Court ordered reargument on whether to reconsider \textit{National League of Cities}.

\section{Judicial Votes}
Reversed, 5--4. Blackmun wrote for Brennan, White, Marshall, and Stevens. Powell (with Burger, Rehnquist, and O'Connor), Rehnquist, and O'Connor dissented.

\section{Holding}
\textit{National League of Cities} is overruled. Immunity based on ``traditional governmental functions'' is unworkable and unsound; state sovereignty is protected chiefly by the national political process, and applying the FLSA to SAMTA violated no constitutional limit.

\section{Analysis of Opinions}
\textbf{Majority:} Lower courts reached irreconcilable results, much as the old test separating governmental from proprietary functions in tax immunity had failed, and judicial line-drawing penalizes state experimentation. The Framers protected the States through the structure of the federal government---the Senate, the Electoral College, and state control of House qualifications---and the States had in fact won vast grants and statutory exemptions. \textbf{Powell:} the Court discarded recent precedent, left Congress the judge of its own power, and ignored how national politics had eroded Wechsler's premises. \textbf{Rehnquist:} the dissenters' principle will again command a majority. \textbf{O'Connor:} the Necessary and Proper Clause incorporates the Tenth Amendment's protection of state autonomy.

\section{Examples: Future Applications}
\textbf{Same-Side:} a generally applicable fleet-electrification rule binds public bus operators. \textbf{Opposite-Side:} ordering state DMVs to enforce that rule against private operators is commandeering (\textit{Printz}). \textbf{Fence-Sitter:} pay and hours rules aimed only at state-court and legislative staff.

\section{Critique}
Critics say the Seventeenth Amendment and nationalized politics have weakened the political safeguards, and Blackmun never defined a process failure. Modern doctrine also asks how Congress regulates rather than how heavily it burdens the States.

\section{Key Quotations}
\begin{quote}
``We therefore now reject, as unsound in principle and unworkable in practice, a rule of state immunity from federal regulation that turns on a judicial appraisal of whether a particular governmental function is `integral' or `traditional.' Any such rule leads to inconsistent results at the same time that it disserves principles of democratic self-governance, and it breeds inconsistency precisely because it is divorced from those principles.''
\end{quote}
\noindent --- Justice Blackmun, for the Court, 469 U.S. at 546--47.

\end{document}
